%\begin{abstract}{Title}{%
%		Speaker}{%
%		University}{%
%		\KLtag}
%Bio
%\end{abstract}

\begin{abstract}{Keynote 1: Interlinking hybrid distribution transformer for hybrid power grids}{%
		Prof. Mariusz Malinowski}{%
		Professor at the Institute of Control and Industrial Electronics, WUT, Poland.
	}{%
		\KLtag}
	\textbf{About the Speaker}:
	Professor Mariusz Malinowski received his Ph.D. and D.Sc. degrees in electrical engineering from the Warsaw University of Technology (WUT), Poland. He has authored over 200 technical papers and seven books. His research focuses on the control and modulation of grid-side converters, multilevel converters, smart grids, and renewable energy systems. He has received numerous awards, including the IEEE IES David Irwin Early Career Award, IEEE IES Bimal Bose Energy Systems Award, and the Polish Prime Minister Award. He is a Fellow of IEEE and a member of the Polish Academy of Sciences.
	
\end{abstract}

\begin{abstract}{Keynote 2: Brushless Field-excitation Systems for Motors in Electric Vehicles}{%
		Prof. Baylon G. Fernandes}{%
		Professor at the Department of Electrical Engineering, IIT Bombay, India.}{%
		\KLtag}
	\textbf{About the Speaker}:
	Professor Baylon G. Fernandes received the B.Tech. degree from the National Institute of Technology Karnataka, Surathkal, India, in 1984, the M.Tech. degree from the Indian Institute of Technology Kharagpur, India, in 1989, and the Ph.D. degree from the Indian Institute of Technology Bombay (IIT Bombay), India, in 1993, all in electrical engineering. Since 1998, he has been with the Department of Electrical Engineering, IIT Bombay, where he is currently a Professor. From 2015 to 2020, he served as the Chairman of the Department. He is a Principal Investigator of the National Centre for Photovoltaic Research and Education and the National Centre of Excellence in Technology for Internal Security at IIT Bombay. His research interests include power electronic converters, motors for electric vehicles, and power electronic interfaces for non-conventional energy sources. Prof. Fernandes is a Senior Member of IEEE and serves as an Associate Editor for IEEE Transactions on Power Electronics.
\end{abstract}

\begin{abstract}{Keynote 3: Transformative Role of Power Electronics in Accelerating Economic Growth of Bharat}{%
		Prof. Rajendra Singh}{%
		Houser Banks Distinguished Professor at Clemson University, USA
	}{%
		\KLtag}
	\textbf{About the Speaker}:
	Professor Rajendra Singh is the Houser Banks Distinguished Professor in the Holcombe Department of Electrical and Computer Engineering and Automotive Engineering at Clemson University (CU). Inspired by Thomas Edison's vision of solar energy, he pursued his Ph.D. dissertation on Silicon Solar Cells during the 1973 energy crisis. Over the past 50 years, he has significantly contributed to the growth of the photovoltaic and semiconductor industries. With expertise in operations, project leadership, R\&D, product commercialization, and startups, Dr. Singh is dedicated to mitigating climate challenges through sustainable green electric power. He has authored or co-authored over 500 publications and edited more than fifteen conference proceedings. He has delivered over 60 keynote addresses and invited talks at national and international conferences. Dr. Singh is the founding technical chair of the IEEE DC Microgrid Conference and currently serves as Chair of the IEEE Power and Energy Society Committee on End-to-End DC Power. He holds six patents, with technologies from his lab licensed for commercialization. Recognized with numerous national and international awards, he was named one of the ten Global Champions of Photovoltaics by Photovoltaics World in 2010. In 2014, he was honored by US President Barack Obama as a White House 'Champion of Change for Solar Deployment.' In 2019, he received the Hind Rattan (Jewel of India) Award. Dr. Singh is a Fellow of IEEE, SPIE, ASM, and AAAS.
\end{abstract}

\begin{abstract}{Keynote 4: Solid State transformer: A Promising Technology Needing Many Breakthroughs}{%
		Prof. Drazen Dujic}{%
		Associate Professor and Head of the Power Electronics Laboratory at EPFL, Switzerland.
	}{%
		\KLtag}
	\textbf{About the Speaker}:
	Drazen Dujic received the Dipl. Ing. and MSc degrees from the University of Novi Sad, Serbia, and the PhD degree from Liverpool John Moores University, UK. He has worked as a Research Assistant at the University of Novi Sad, a Research Associate at Liverpool John Moores University, and in various roles at ABB Switzerland Ltd. Since 2014, he has been with EPFL. His research focuses on high-power electronic systems and high-performance drives. He has received multiple awards, including the Istvan Nagy Award (2024), the EPE Outstanding Service Award (2018), and the Isao Takahashi Power Electronics Award (2014). He is an IEEE Fellow.
\end{abstract}



\begin{abstract}{Keynote 5: Magnetic energy harvesting from overhead power lines}{%
		Alon Kuperman}{%
		Ben-Gurion University of the Negev, Israel}{%
		\KLtag}
	\textbf{About the Speaker}: Alon Kuperman (Senior Member, IEEE) is a Full Professor with the School of Electrical and Computer Engineering, Ben-Gurion University of the Negev, Beer-Sheva, Israel and the Director of Applied Energy Laboratory. He holds multiple international patents and owns an independent consultancy services company. He is also a co-founder of two start-up companies focusing on wireless power transfer. His research interests include all aspects of energy conversion and applied control.
\end{abstract}


